

\subsubsection{Semantic-Oriented Transformer Models}

While the two-stage model is designed to explicitly capture low-entropy
editorial and structural patterns, it is inherently limited in modeling
deeper semantic relations.
To complement this inductive bias, we employ transformer-based language
models, which encode contextual and semantic dependencies through
self-attention mechanisms.

We consider two architectures: RoBERTa and DeBERTa.
RoBERTa serves as a strong and widely adopted baseline for contextual text
representations, particularly robust in noisy and HTML-rich settings.
DeBERTa introduces disentangled attention mechanisms that separately model
content and positional information, enabling finer-grained semantic
discrimination.

Both architectures are fine-tuned on the task and used as high-capacity
semantic classifiers, operating independently from the rule-based stage.
Their predictions are subsequently analyzed and combined through overlap
and ensemble strategies, allowing us to exploit complementarity between
editorial-driven and semantic-driven decision processes.
